\documentclass[11pt,a4paper]{article}
\usepackage[margin=25mm,headheight=14pt]{geometry}
\usepackage[T1]{fontenc}
\usepackage[utf8]{inputenc}
\usepackage{lmodern,amsmath,amssymb}
\usepackage{xcolor,graphicx,booktabs,longtable,array}
\usepackage{microtype,enumitem,fancyhdr,titlesec}
\usepackage{hyperref}
\definecolor{accent}{HTML}{163B55}
\hypersetup{colorlinks=true,linkcolor=accent,urlcolor=accent,
 pdftitle={The sharp upper constant for the eigenvalue--torsion product on convex domains},pdfauthor={Matthew J. Colbrook},
 pdfsubject={AIM 1028: mathematical result and proof}}
\urlstyle{same}
\Urlmuskip=0mu plus 1mu
\setlength{\parindent}{0pt}
\setlength{\parskip}{0.35em}
\setlength{\emergencystretch}{3em}
\setlist{itemsep=0.25em,topsep=0.3em}
\providecommand{\tightlist}{\setlength{\itemsep}{0pt}\setlength{\parskip}{0pt}}
\titleformat{\section}{\normalfont\large\bfseries\color{accent}}{\thesection}{0.7em}{}
\titleformat{\subsection}{\normalfont\normalsize\bfseries}{\thesubsection}{0.6em}{}
\pagestyle{fancy}
\fancyhf{}
\fancyhead[L]{\small\sffamily AIM 1028}
\fancyhead[R]{\small\sffamily Matthew J. Colbrook}
\fancyfoot[C]{\thepage}
\renewcommand{\headrulewidth}{0.3pt}
\setcounter{secnumdepth}{0}
\allowdisplaybreaks[2]
\title{\textbf{The sharp upper constant for the eigenvalue--torsion product on convex domains}}
\author{Matthew J. Colbrook\\[3pt]
{\small Department of Applied Mathematics and Theoretical Physics}\\
{\small University of Cambridge, United Kingdom}\\
{\small\href{mailto:mjc249@maths.cam.ac.uk}{\texttt{mjc249@maths.cam.ac.uk}}}}
\date{3 November 2026}
\begin{document}
\maketitle
\begin{abstract}
We establish the sharp upper constant $\pi^2/12$ for the normalized first-eigenvalue--torsional-rigidity product on bounded convex domains in every dimension. The proof explicitly extends the one-dimensional mechanism in Deng's work using a dimension-independent torsion--perimeter inequality. The lower-constant question in AIM 1028 is not settled.
\end{abstract}
\noindent\textbf{Repository target.}
This manuscript concerns AIM problem 1028 at revision
\texttt{3fa50884d129}. Its scope is stated below.
\medskip

\hypertarget{1028-status-and-attribution}{%
\section{Status and attribution}\label{1028-status-and-attribution}}

The \textbf{upper half} of problem 1028 follows in every dimension from
the one-dimensional mechanism in Section 5 of Zikang Deng, \emph{A Proof
of a Conjecture on a Pólya Functional}, arXiv:2609.13221v1 (2026),
together with the dimension-independent Pólya torsion-perimeter
inequality. The argument below checks the extension explicitly and
supplies a direct proof of that geometric input.

\textbf{Theorem.} Let \(d \ge 2\) and let \(\Omega\) be a nonempty
bounded open convex subset of \(\mathbb R^d\). Let \(u\) solve
\(-\Delta u=1\) with zero Dirichlet boundary values, let
\(T=\int _\Omega u\), and let \(\lambda _1\) be the first Dirichlet
eigenvalue. Then

\[
 \frac{\lambda_1(\Omega)T(\Omega)}{|\Omega|}<\frac{\pi^2}{12}.
\]

Moreover,

\[
 \sup_{\Omega\subset\mathbb R^d\ \mathrm{bounded,open,convex}}
 \frac{\lambda_1(\Omega)T(\Omega)}{|\Omega|}=\frac{\pi^2}{12}.
\]

The proposed \textbf{lower} bound in problem 1028 is not established
here. Thus this is a complete resolution of one sharp constant, not a
complete resolution of the combined catalogue entry. No claim of novelty
beyond the checked extension of Deng's argument is made.

\hypertarget{1028-a-geometric-input-with-no-dimension-restriction}{%
\section{1. A geometric input with no dimension
restriction}\label{1028-a-geometric-input-with-no-dimension-restriction}}

For a bounded convex domain \(D\) in \(\mathbb R^d\), write \(V=|D|\),
\(P=\mathcal H^{d-1}(\partial D)\), and
\(\rho(x)=\operatorname{dist}(x,\partial D)\). Let

\[
 D_s=\{\rho>s\},\qquad V(s)=|D_s|,\qquad P(s)=P(D_s),
\]

for \(0<s<r\), where \(r\) is the inradius. The coarea formula gives
\(V'(s)=-P(s)\) almost everywhere. Convexity and monotonicity of surface
area for nested convex bodies give \(P(s)\le P(0)=P\). In fact the
inequality is strict for \(s>0\) while \(D_s\) has interior. For
example, \(D_s+sB\) lies in \(D\), and the surface area of \(D_s+sB\) is
strictly greater than that of \(D_s\) when \(d\ge 2\).

Use the distance-dependent trial

\[
 v(x)=\int_0^{\rho(x)}\frac{V(s)}{P(s)}\,ds.
\]

This belongs to \(H^1_0(D)\): since
\(V(s)=\int _s^r P(t)dt \le (r-s)P(s)\), its defining derivative is
bounded. Fubini and coarea give

\[
 \int_D v=\int_D|\nabla v|^2
  =\int_0^r\frac{V(s)^2}{P(s)}\,ds.
\]

The variational characterization of torsional rigidity consequently
yields

\[
 T(D)\ge\int_0^r\frac{V(s)^2}{P(s)}\,ds
       =\int_0^{V}\frac{v^2}{P(s(v))^2}\,dv
       >\frac{V^3}{3P^2}.
 \tag{1}
\]

The weak version suffices for the sharp supremum; strictness gives
nonattainment. This is the all-dimensional Pólya torsion-perimeter
inequality. It also appears, more generally for the p-Laplacian, in
Amato--Masiello--Paoli--Sannipoli (2023), Introduction equation (8) and
Theorem 1.1 with \(f=1,p=2\).

\hypertarget{1028-torsion-superlevels-and-an-ordinary-differential-inequality}{%
\section{2. Torsion superlevels and an ordinary differential
inequality}\label{1028-torsion-superlevels-and-an-ordinary-differential-inequality}}

Put

\[
 A=|\Omega|,\quad M=\max_\Omega u,\quad U=u/M,
\]

and define for \(0\le t\le 1\)

\[
 S(t)=\frac{|\{U>t\}|}{A},\qquad
 I(t)=\int_t^1 S(s)\,ds,\qquad
 \Phi=I(0)=\frac{T}{AM}.
 \tag{2}
\]

The power-concavity theorem of Kennington states that \(\sqrt{u}\) is
concave on a convex domain in every dimension. In particular, every
superlevel domain \(\Omega _t={U>t}\) is convex.

For completeness, the level sets with \(0<t<1\) are regular. If two
points maximized \(u\), concavity of \(\sqrt{u}\) would make it constant
on the segment between them. Interior analyticity then makes its
restriction constant along the entire maximal chord containing the
segment, contradicting zero boundary values at the chord endpoints.
Hence the maximum is unique. A critical point of \(u\) is a maximum,
because a differentiable concave function lies below its tangent
hyperplane. Therefore \(\nabla u\) is nonzero away from this unique
point. This argument is valid in every dimension.

The function \(u-Mt\) is the torsion function of \(\Omega _t\). Thus

\[
 |\Omega_t|=AS(t),\qquad T(\Omega_t)=AM I(t).
 \tag{3}
\]

Flux and coarea, with \(P_t=P(\Omega _t)\), give

\[
 \int_{\partial\Omega_t}|\nabla u|\,dS=AS(t),\qquad
 \int_{\partial\Omega_t}|\nabla u|^{-1}\,dS
       =-\frac{A}{M}S'(t).
\]

Cauchy--Schwarz implies

\[
 P_t^2\le\frac{A^2}{M}S(t)(-S'(t)).
\]

Combining this with (1) applied to \(\Omega _t\) gives

\[
 I(t)(-S'(t))>\frac{S(t)^2}{3},
 \qquad\text{equivalently}\qquad
 I(t)I''(t)>\frac{I'(t)^2}{3}.
 \tag{4}
\]

All powers of \(A\) and \(M\) cancel; there is no dependence on \(d\)
here. Therefore

\[
 \big(I^{2/3}\big)''
  =\frac23 I^{-4/3}\left(II''-\frac13 I'^2\right)>0.
 \tag{5}
\]

The standard approximation by smooth convex outer domains can be used
throughout if desired. Alternatively, all level sets above are compact
smooth hypersurfaces in the interior, and the identities hold directly
for a bounded convex domain.

\hypertarget{1028-the-second-moment-factor}{%
\section{3. The second-moment factor}\label{1028-the-second-moment-factor}}

At the left endpoint, \(I(0)=\Phi\) and \(I'(0+)=-1\); at the right
endpoint, \(I(1)=0\). The tangent line at zero to the strictly convex
function \(I^{2/3}\) gives

\[
 I(t)\ge\Phi\left(1-\frac{2t}{3\Phi}\right)_+^{3/2},
 \qquad \Phi\le\frac23.
\]

The inequality is strict at every positive \(t\) before the positive
part vanishes. Integrating,

\[
 \int_0^1 I(t)\,dt>\frac{3\Phi^2}{5}.
\]

Writing \(m_k=A^{-1}\int _\Omega U^k\), layer cake gives

\[
 m_1=\Phi,\qquad
 m_2=2\int_0^1tS(t)\,dt=2\int_0^1 I(t)\,dt.
\]

Consequently,

\[
 \boxed{\frac{T^2}{A\int_\Omega u^2}
       =\frac{\Phi^2}{m_2}<\frac56.}
 \tag{6}
\]

\hypertarget{1028-a-weighted-coordinate}{%
\section{4. A weighted coordinate}\label{1028-a-weighted-coordinate}}

Introduce the decreasing coordinate

\[
 y(t)=\left(\frac{I(t)}{\Phi}\right)^{1/3},
 \qquad t=t(y),
 \qquad
 a(y)=\frac{S(t(y))}{3\Phi y}.
 \tag{7}
\]

Then

\[
 \frac{dt}{dy}=-\frac{y}{a(y)},\qquad
 S(t(y))=3\Phi y a(y),\qquad
 \int_0^1\frac{y}{a(y)}\,dy=1.
 \tag{8}
\]

Since \((I^{2/3})'=-2 \Phi ^{2/3} a(y(t))\), convexity in (5), together
with reversal of orientation between \(y\) and \(t\), says that \(a\) is
nondecreasing. The endpoints satisfy

\[
 a(1)=\frac1{3\Phi},\qquad
 y a(y)\longrightarrow0\quad (y\downarrow0).
\]

The change of Stieltjes measure is

\[
 \int_{[0,1]}\psi(t)(-dS(t))
 =3\Phi\int_{[0,1]}\psi(t(y))\,d(ya(y)).
 \tag{9}
\]

For every Lipschitz \(H\) with \(H(0)=0\), the trial function \(H(U)\)
and coarea give

\[
 \lambda_1M
 \le\frac{\int_0^1 S(t)H'(t)^2\,dt}
          {\int_{[0,1]}H(t)^2(-dS(t))}.
\]

Putting \(f(y)=H(t(y))\), equations (8)--(9) reduce this to

\[
 \lambda_1M\le
 \frac{\int_0^1 a(y)^2 f'(y)^2\,dy}
      {\int_{[0,1]}f(y)^2\,d(ya(y))}.
 \tag{10}
\]

Another change of variables and an integration by parts give

\[
 \frac{m_2}{m_1}
 =2\int_0^1\frac{y^4}{a(y)}\,dy=2\kappa,
 \qquad
 \kappa:=\int_0^1\frac{y^4}{a(y)}\,dy.
 \tag{11}
\]

\hypertarget{1028-the-sharp-weighted-inequality}{%
\section{5. The sharp weighted
inequality}\label{1028-the-sharp-weighted-inequality}}

Let

\[
 h(y)=\sin(\pi y/2),\qquad
 f(y)=\int_y^1\frac{h(s)}{a(s)}\,ds.
 \tag{12}
\]

This is admissible in (10). Indeed \(h(y)\le \pi y/2\) and (8) give
\(f(0)\le \pi /2\). In the original coordinate,

\[
 H'(t)=\frac{h(y)}{y}\le\frac\pi2,
 \qquad H(0)=f(1)=0.
\]

The numerator in (10) is exactly

\[
 \int_0^1a^2 f'^2=\int_0^1h^2=\frac12.
 \tag{13}
\]

For the denominator, Stieltjes integration by parts and Fubini yield

\[
 D:=\int_{[0,1]}f^2\,d(ya)
   =2\int_0^1\frac{h(s)}{a(s)}
                  \left(\int_0^s y h(y)\,dy\right)ds.
 \tag{14}
\]

There is no endpoint term because \(f(1)=0\), \(f(0)<\infty\), and
\(ya(y)\to 0\).

Define

\[
 r(s)=2h(s)\int_0^s yh(y)\,dy-\frac{10}{\pi^2}s^4,
\]

and compute its primitive explicitly:

\[
 R(x):=\int_0^x r(s)\,ds
   =\frac{2}{\pi^3}
     \left[\pi x\{2+\cos(\pi x)-x^4\}-3\sin(\pi x)\right].
 \tag{15}
\]

We need \(R\ge 0\) on \([0,1]\), with \(R(0)=R(1)=0\). Here is an
elementary verification. Set

\[
 q(z)=\frac{z(2+\cos z)-3\sin z}{z^5},\qquad 0<z\le\pi.
\]

Then

\[
 z^6 q'(z)=-G(z),\quad
 G(z)=z(8+7\cos z)-(15-z^2)\sin z.
\]

Direct differentiation gives

\[
 G'''(z)=z(\sin z-z\cos z)\ge0,
 \qquad G(0)=G'(0)=G''(0)=0,
\]

because \(\sin z-z \cos z=\int _0^z s \sin s ds\ge 0\) on \([0,\pi ]\).
Thus \(G\ge 0\), so \(q\) is nonincreasing and

\[
 q(z)\ge q(\pi)=\pi^{-4}.
\]

With \(z=\pi x\), this is exactly \(R(x)\ge 0\) in (15).

Now \(w=1/a\) is nonincreasing. Integration by parts gives

\[
 \int_0^1\frac{r(s)}{a(s)}\,ds
  =-\int_{[0,1]}R(s)\,dw(s)\ge0.
 \tag{16}
\]

To justify the possibly singular lower endpoint, monotonicity and (8)
imply

\[
 \frac{y^2w(y)}2\le\int_0^y sw(s)\,ds\longrightarrow0;
\]

also \(R(y)=O(y^5)\). Hence \(R(y)w(y)\to 0\). The upper endpoint
vanishes because \(R(1)=0\) and \(w(1)=3\Phi\) is finite.

Equations (14)--(16) give

\[
 D\ge\frac{10}{\pi^2}\kappa.
\]

Using (10), (11), and (13), we conclude

\[
 \boxed{\frac{\lambda_1\int_\Omega u^2}{T}
       =2\kappa\lambda_1M
       \le\frac{\kappa}{D}\le\frac{\pi^2}{10}.}
 \tag{17}
\]

Multiplying (6) and (17) proves

\[
 \frac{\lambda_1T}{A}
  =\frac{T^2}{A\int_\Omega u^2}
      \frac{\lambda_1\int_\Omega u^2}{T}
  <\frac56\frac{\pi^2}{10}=\frac{\pi^2}{12}.
\]

\hypertarget{1028-sharpness-in-every-dimension}{%
\section{6. Sharpness in every
dimension}\label{1028-sharpness-in-every-dimension}}

Take the rectangular boxes

\[
 \Omega_L=(-L/2,L/2)^{d-1}\times(0,1),\qquad L\to\infty.
\]

Separation of variables gives

\[
 \lambda_1(\Omega_L)=\pi^2\left(1+\frac{d-1}{L^2}\right).
\]

Let \(u_L\) be their torsion functions and \(w(y)=y(1-y)/2\), the
infinite-slab torsion function. By comparison, \(0\le u_L\le w\). The
nonnegative difference \(w-u_L\) is harmonic. Choose

\[
 C=\max_{0\le y\le1}\frac{w(y)}{\sin(\pi y)}<\infty.
\]

The harmonic barrier

\[
 B_L(x,y)=C\sin(\pi y)
     \sum_{j=1}^{d-1}\frac{\cosh(\pi x_j)}{\cosh(\pi L/2)}
\]

dominates \(w-u_L\) on every side of the box. Its integral is
\(O(L^{d-2})\), whereas the volume is \(L^{d-1}\). Consequently

\[
 \frac{T(\Omega_L)}{|\Omega_L|}
 =\int_0^1w(y)\,dy+O(L^{-1})
 =\frac1{12}+O(L^{-1}).
\]

It follows that
\(\lambda _1(\Omega _L)T(\Omega _L)/|\Omega _L| \to \pi ^2/12\).

\hypertarget{1028-references-and-source-scope}{%
\section{References and source
scope}\label{1028-references-and-source-scope}}

\begin{enumerate}
\def\labelenumi{\arabic{enumi}.}
\tightlist
\item
  Z. Deng, \emph{A Proof of a Conjecture on a Pólya Functional},
  arXiv:2609.13221v1 (2026), Section 5.
  \url{https://arxiv.org/abs/2609.13221}. The source states the planar
  theorem; the calculations above explicitly check the upper-bound
  mechanism in arbitrary dimension.
\item
  V. Amato, A. L. Masiello, G. Paoli, R. Sannipoli, \emph{Sharp and
  quantitative estimates for the p-Torsion of convex sets}, Nonlinear
  Differential Equations and Applications NoDEA \textbf{30} (2023),
  Article 12. \url{https://doi.org/10.1007/s00030-022-00822-y}. Theorem
  1.1 and Introduction equation (8) are all-dimensional; normalization
  at \(f=1,p=2\) is \(T=\int u\).
\item
  A. U. Kennington, \emph{Power concavity and boundary value problems},
  Indiana University Mathematics Journal \textbf{34} (1985), 687--704.
  Theorem~4.1(ii), pp.~695--696, gives the torsion square-root
  concavity theorem in arbitrary dimension.
  \url{https://doi.org/10.1512/iumj.1985.34.34036}.
  \url{https://math.jhu.edu/~js/Math646/kennington.pdf}.
\item
  M. van den Berg, G. Buttazzo, A. Pratelli, \emph{On the relations between
  principal eigenvalue and torsional rigidity}, Communications in
  Contemporary Mathematics \textbf{23} (2021), 2050093, Conjecture 4.2.
  \url{https://arxiv.org/abs/1910.14593}.
\end{enumerate}

The lower constant \(\pi ^2/[2(d+1)(d+2)]\) remains unproved by this
argument. The planar Airy-stress step used for that constant in Deng's
paper does not directly generalize to higher dimensions.
\appendix
\section{A limitation of a lower-bound proof route}
This appendix records a limitation of one attempted route to the lower half of AIM 1028. It does not claim a counterexample to the conjecture or an improved general lower bound.

Let \(u\) be the torsion function of a bounded convex domain \(\Omega\subset\mathbb R^d\), and write
\[
M=\max_\Omega u,\qquad A=|\Omega|,\qquad T=\int_\Omega u.
\]
The proposed estimate would follow from the two inequalities
\[
\lambda_1 M\ge\frac{\pi^2}{8},\qquad
\frac{T}{AM}\ge\frac{4}{(d+1)(d+2)}.
\]
The first estimate has a dimension-independent proof. The second, which is sharp on thin cone families if true, was not established here for \(d\ge3\).

\subsection{1. What the verified general result supplies}

Della Pietra, Gavitone and Guarino Lo Bianco, \emph{On functionals involving the torsional rigidity related to some classes of nonlinear operators}, J. Differential Equations \textbf{265} (2018), 6424--6442, DOI \href{https://doi.org/10.1016/j.jde.2018.07.030}{10.1016/j.jde.2018.07.030}, \href{https://arxiv.org/abs/1705.03330}{arXiv:1705.03330}, Theorems 1.1(i) and 4.2, prove
\[
\frac{T_p}{A M}\ge \frac{q}{d^{q-1}(d+q)},\qquad q=\frac{p}{p-1}.
\]
At \(p=q=2\), this is
\[
\frac{T}{AM}\ge\frac{2}{d(d+2)}.
\]
Their Remark 4.6 explicitly displays the factor \(2\); a sentence in the introduction drops it, so the theorem rather than that sentence should be used. Briani and Bucur, \emph{Mean-to-max ratio of the torsion function and honeycomb structures}, Calc. Var. \textbf{62} (2023), 198, DOI \href{https://doi.org/10.1007/s00526-023-02530-6}{10.1007/s00526-023-02530-6}, Remark 5.2, discusses the still unknown sharp convex lower bound and the same \(p\)-torsion estimate.

Combining the verified bound with Payne's estimate gives
\[
\frac{\lambda_1T}{A}\ge\frac{\pi^2}{4d(d+2)}.
\]
This is below the conjectured constant \(\pi^2/[2(d+1)(d+2)]\). The needed improvement factor is \(2d/(d+1)\).

For completeness, the first factor follows from the dimension-independent gradient bound
\[
|\nabla u|^2+2u\le2M
\]
(Theorem 3.2 of the same 2018 paper, with \(p=2\)). For a nonsmooth convex domain, smooth convex approximation and interior
\(C^1\) convergence, with maxima controlled by domain monotonicity and
scaling, pass this gradient inequality to the limit. Put
\[
p(t)=\cos\!\left(\frac\pi2\sqrt{1-t/M}\right),\qquad w=p(u).
\]
Then \(p''>0\), \(w\in H^1_0(\Omega)\), and
\[
-\Delta w-\frac{\pi^2}{8M}w
=p''(u)\bigl(2(M-u)-|\nabla u|^2\bigr)\ge0.
\]
Pairing with the positive first Dirichlet eigenfunction proves \(\lambda_1M\ge\pi^2/8\). This is also the proof in Proposition 2.2 of Deng's planar paper, but the calculation and gradient estimate are valid in every dimension.

\subsection{2. A rigorous obstruction to a proof using only the two usual shape inequalities}

The combination of square-root concavity and the gradient inequality does \textbf{not}, by itself, imply the desired torsion mean-to-maximum bound. This can be shown explicitly.

For \(L>\sqrt2\), define a radial function on \(B_L\subset\mathbb R^d\) by
\[
v_L(x)=v_L(r)=
\begin{cases}
1-r^2/2,&0\le r\le 2/L,\\[2mm]
\displaystyle\frac{(L-r)^2}{L^2-2},&2/L\le r\le L,
\end{cases}
\qquad r=|x|.
\]
The values and first derivatives match at \(r=2/L\): they are respectively \(1-2/L^2\) and \(-2/L\). Thus \(v_L\) is \(C^1\), belongs to \(H^1_0(B_L)\), is positive inside the ball, and has maximum one.

Its square root is concave. On the inner interval it equals \(\sqrt{1-r^2/2}\), a concave decreasing function of \(r\); on the outer interval it is the affine function \((L-r)/\sqrt{L^2-2}\). The radial derivatives match at the interface, and the tangential Hessian eigenvalues have the sign of the nonpositive radial derivative.

Moreover,
\[
|\nabla v_L|^2+2v_L\le2=2\max v_L.
\]
There is equality on the inner interval. On the outer interval, writing \(R^2=L^2-2\) and \(w=\sqrt{v_L}\), the left side is \(w^2(2+4/R^2)\), and
\[
w^2\le v_L(2/L)=\frac{R^2}{R^2+2},
\]
which proves the estimate.

Nevertheless, dominated convergence after the substitution \(r=Ls\) gives
\[
\frac{1}{|B_L|}\int_{B_L}v_L
=d\int_0^1v_L(Ls)s^{d-1}\,ds
\longrightarrow
d\int_0^1(1-s)^2s^{d-1}\,ds
=\frac{2}{(d+1)(d+2)}.
\]
This is half the proposed efficiency lower bound. In particular, for all sufficiently large \(L\), these functions violate the proposed efficiency inequality while satisfying both square-root concavity and the gradient estimate.

These functions are \textbf{not torsion functions}: on the inner ball, \(-\Delta v_L=d\), while on the outer annulus,
\[
-\Delta v_L
=\frac{2}{L^2-2}\left(\frac{(d-1)(L-r)}{r}-1\right),
\]
which is nonconstant and negative near the outer boundary. Therefore this construction does not disprove the conjecture. It pinpoints the limitation of the proposed proof route: the torsion equation must supply additional information beyond its square-root concavity and the usual gradient bound.

\subsection{3. Status}

The all-dimensional sharp upper bound is proved in the preceding sections. The lower bound in dimensions \(d\ge3\) remains unresolved in this work. Deng's planar argument obtains the stronger efficiency estimate through a planar Airy-stress construction and mixed-area geometry; merely retaining its dimension-independent concavity and gradient consequences cannot complete the higher-dimensional lower estimate.
\end{document}
